% =====================================================================
%  CAPÍTULO 6 — IMPLEMENTAÇÃO
% =====================================================================
\chapter{Implementação}
\label{cap:implementacao}

Este capítulo descreve como os requisitos e o \emph{design} foram
traduzidos em código. Organiza-se pelos quatro módulos do \emph{pipeline}
(M1--M4), seguidos da infraestrutura do \emph{backend}, do \emph{frontend}
e das decisões de segurança. Ao longo do texto incluem-se excertos de
código \emph{críticos} — não a totalidade da implementação, mas os pontos
que melhor ilustram as decisões técnicas.

\section{Módulo M1 --- Ingestão Multimodal}
\label{sec:impl_m1}

O M1 é o ponto de entrada de todo o material. Recebe os \emph{bytes} de um
ficheiro e orquestra uma cadeia de processamento: \emph{security scan}
$\rightarrow$ \textsc{exif} $\rightarrow$ carregamento para o \emph{Storage}
$\rightarrow$ análise visual $\rightarrow$ \textsc{ocr}, persistindo um
registo \texttt{MediaFile}.

O processamento decorre sobre um ficheiro temporário (as bibliotecas de
\textsc{exif}/\textsc{ocr}/visão esperam um caminho de sistema), removido no
final. Por segurança, os \emph{bytes} só chegam ao \emph{Storage} \emph{após}
um \emph{scan} bem-sucedido. As etapas são:

\begin{description}[leftmargin=*]
  \item[\emph{Security scan}.] Cálculo de \emph{checksum} MD5 e \emph{scan}
  antivírus com o \textbf{ClamAV} (quando presente no sistema; na sua
  ausência, o passo degrada graciosamente e regista um aviso). Perante a
  deteção de uma ameaça, o registo é marcado \texttt{FAILED} e o
  \emph{pipeline} interrompe-se --- o conteúdo nunca chega ao \emph{Storage}.
  \item[\textsc{exif}.] Data de captura, coordenadas GPS, marca e modelo
  da câmara, extraídos com a biblioteca \textbf{exifread}~\cite{exifread} (puro Python, sem
  dependências de sistema, funcionando de forma idêntica em ambiente local
  e na nuvem) e preservando o \textsc{exif} em bruto. As coordenadas, no
  formato \emph{graus/minutos/segundos}, são convertidas para decimal com o
  sinal correto segundo a referência (N/S, E/W).
  \item[Análise visual.] O \textbf{Gemini Vision} produz descrição da
  cena, contagem de pessoas, cenário, emoção dominante, \emph{tags} e uma
  \emph{sugestão narrativa} — a matéria-prima factual do M3.
  \item[\textsc{ocr}.] Para documentos, o \textbf{Tesseract}~\cite{smith2007tesseract} extrai
  texto, enriquecendo o contexto factual.
\end{description}

\textbf{Análise diferida.} A análise visual é o passo mais lento (uma
chamada ao Gemini, da ordem dos segundos). Para não bloquear o carregamento
--- sobretudo em lotes ---, o M1 corre no pedido HTTP apenas a parte rápida
(\emph{scan}, \textsc{exif} e \emph{Storage}) e \emph{difere} a IA
(Gemini/\textsc{ocr}) para segundo plano, no executor \emph{in-process}
(\Cref{sec:impl_backend}). A fotografia fica logo visível; o registo
permanece em \texttt{PROCESSING} e a descrição preenche-se no fim,
sinalizada por um indicador ``a analisar'' na interface, que faz
\emph{polling} enquanto houver registos por concluir.

\textbf{Re-análise sob procura.} Como o executor de segundo plano pode ser
terminado pela plataforma gratuita antes de concluir (deixando fotografias
presas em \texttt{PROCESSING}), um \emph{endpoint} de \textbf{re-análise}
--- e um botão na Biblioteca --- reprocessa de forma síncrona as que ficaram
sem descrição, marcando-as \texttt{COMPLETED} e desbloqueando a linha
temporal e o vídeo. É \emph{consciente da quota} (indica quantas serão
analisadas antes de avançar) e processa em \textbf{pequenos lotes} por
pedido, para não esgotar a memória da instância gratuita.

A árvore genealógica é tratada por um \emph{parser} \textbf{GEDCOM}~\cite{gedcom}
dedicado (o maior módulo, $\approx$430 linhas), que converte o ficheiro em
registos \texttt{Person} e nas respetivas relações, com suporte a
etiquetagem de famílias (permitindo distinguir várias árvores na mesma
conta).

\section{Módulo M2 --- Organização Temporal}
\label{sec:impl_m2}

O M2 transforma factos dispersos numa estrutura temporal e relacional
coerente, através de dois componentes.

O \textbf{\emph{date resolver}} normaliza datas heterogéneas e
frequentemente incompletas, produzindo uma chave de ordenação consistente
— essencial, dado que o material de arquivo raramente tem datas exatas.
Trata os qualificadores do formato GEDCOM (\texttt{ABT} ``cerca de'',
\texttt{BEF} ``antes de'', \texttt{AFT} ``depois de'') e os casos de apenas
mês/ano ou só ano, e rejeita datas implausíveis (anteriores a 1800, o
\emph{epoch} Unix de 1970 ou no futuro), evitando que ruído contamine a
linha temporal. Os casamentos importados do GEDCOM geram automaticamente
\texttt{TimelineEvent}s datados, que passam a fazer parte do contexto
factual disponível ao M3.

O coração do M2 é o \textbf{grafo de parentesco} (\texttt{FamilyGraph}),
construído sobre o \emph{NetworkX}~\cite{networkx} como grafo dirigido: os
nós são pessoas e as arestas representam relações (\emph{parent},
\emph{child}, \emph{spouse}, \emph{sibling}). A classe produz
\emph{resumos narrativos} do parentesco (por exemplo, \emph{``João (pai de
Maria, marido de Ana)''}), consumidos diretamente pelo M3 para dar ao
\textsc{llm} o contexto de \emph{quem é quem}.

\textbf{Relações na base de dados e edição manual.} As relações passaram a
ser persistidas numa tabela \texttt{relationships} --- a fonte de verdade
durável ---, em vez de um grafo JSON em disco que, sendo efémero na nuvem,
deixava o contexto familiar vazio após cada reinício da instância (uma
fonte subtil de narrativas sem relações). Isto sustenta, na página
\emph{Família}, uma \textbf{vista interativa da árvore} (\emph{React~Flow},
\emph{layout} por gerações, casais lado a lado, nós coloridos por sexo e
ligações desenhadas com um \textbf{barramento} partilhado --- descida do
casal, barra de irmãos e ligação curta a cada filho ---, evitando linhas
sobrepostas) e um \textbf{editor manual} de pessoas e ligações
(\emph{pai}/\emph{mãe}/\emph{filho}/\emph{irmão}/\emph{casado(a)}), com um
assistente de \textbf{criação rápida de árvore} (eu, pais, avós, irmãos).
Após cada edição, o \texttt{FamilyGraph} é reconstruído da base de dados,
isolado \emph{por utilizador}, mantendo o M3 sincronizado.

As \textbf{notas} de cada pessoa entram no contexto do M3, influenciando a
narrativa. Cada fotografia pode ser \textbf{etiquetada com quem nela
aparece} (\texttt{media\_files.person\_ids}), ligando \emph{rostos a nomes}
--- a ponte entre a árvore e a narrativa ---, e cada pessoa pode ter uma
\textbf{foto de perfil} (\texttt{persons.photo\_media\_id}), usada como
\emph{avatar} na lista e nos nós da árvore. A árvore pode ainda ser
\textbf{exportada para GEDCOM~5.5.1} (re-importável noutras plataformas),
demonstrando \emph{interoperabilidade} com o formato standard.

\textbf{Isolamento de múltiplas árvores.} Importar vários ficheiros GEDCOM
levantava um problema subtil: as exportações genealógicas reiniciam quase
sempre os identificadores em \texttt{@I1@} por ficheiro, pelo que a
deduplicação por \texttt{(user\_id, gedcom\_id)} fazia a segunda família
\emph{sobrescrever} a primeira, fundindo árvores sem relação. A solução foi
\textbf{prefixar o \texttt{gedcom\_id} persistido com a etiqueta da família}
(\texttt{"<família>::<id>"}; sem etiqueta, um \emph{id} de lote único), de
modo que reimportar a \emph{mesma} família continua idempotente mas famílias
distintas nunca colidem. Sobre este mecanismo, cada \textbf{projeto} mantém
a sua família isolada e, dentro dele, cada GEDCOM importado fica num
\textbf{sub-grupo próprio} (nomeado pelo ficheiro), alternável por
\emph{chips} — uma vista por sub-família ou todas em conjunto (parâmetro
\texttt{group} na \textsc{api}, que casa a etiqueta do projeto e as suas
sub-etiquetas).

\section{Módulo M3 --- Geração Narrativa}
\label{sec:impl_m3}

O M3 é o módulo de IA generativa central. Combina três componentes: um
\emph{cliente LLM} híbrido, o subsistema \emph{RAG} e um conjunto de
\emph{templates}.

\subsection{Cliente LLM híbrido e \emph{fallback}}

O \texttt{LLMClient} implementa uma \textbf{cadeia de \emph{fallback} de
três níveis} para o texto: o motor primário é o \emph{Groq}
(\emph{Llama~3.3~70B}, camada gratuita rápida); na sua ausência ---
tipicamente sem Internet --- recorre-se ao \emph{Ollama} local
(Llama~3.2~3B), \emph{offline e privado}; e só se ambos falharem entra o
\emph{Gemini~Flash}. A ordem é \textbf{deliberada}: o Llama~3.2~3B local, em
\textsc{cpu}, é demasiado lento para uso interativo, pelo que o Groq (70B na
nuvem) dá melhor qualidade e resposta quase imediata; e a quota do Gemini
fica \textbf{reservada para a visão} (M1), que se mantém \emph{Gemini-only}
por o Groq não oferecer visão --- repartindo texto e visão por dois
serviços gratuitos sem competirem pela mesma quota. Como mostra o
\Cref{lst:llm}, se um nível falhar \emph{durante} a geração há
\emph{fallback} em tempo de execução; só quando \emph{todos} falham é
lançada uma exceção que marca a tarefa como falhada, evitando guardar uma
narrativa falsa.

\subsection{Retrieval-Augmented Generation com degradação graciosa}

O \texttt{RAGSystem} indexa os factos extraídos pelo M1/M2 numa base
vetorial (ChromaDB~\cite{chromadb}), com o \texttt{user\_id} na
\emph{metadata} de cada documento; a pesquisa filtra sempre por esse campo
(isolamento multi-inquilino). Quando a biblioteca é grande e o utilizador
indica um foco, o método \texttt{search\_media\_ids} injeta no \emph{prompt}
apenas o subconjunto de fotografias mais relevante ao tema (por
similaridade vetorial), mantendo o contexto focado, dentro da janela do
modelo e menos sujeito a alucinação.

A decisão de engenharia mais interessante é a \textbf{degradação graciosa}.
Quando o ChromaDB não está instalado --- por exemplo, no \emph{deploy}
gratuito do Render, onde o custo de \emph{build} é incomportável ---, uma
importação condicional deteta a ausência e comuta o \texttt{RAGSystem} para
um \emph{stub}: a indexação torna-se inócua e a pesquisa devolve listas
vazias, passando o M3 a usar a totalidade dos factos. O sistema mantém-se
\emph{funcional}, perdendo apenas o passo de \emph{retrieval}.

\subsection{\emph{Templates} e calibração para português europeu}

Existem \textbf{sete} \emph{templates}: seis temáticos (\emph{Memória
Familiar}, \emph{Momento Capturado}, \emph{História de Amor},
\emph{Aventura Familiar}, \emph{Nova Vida}, \emph{Momento de Alegria}) e
um \emph{livre} (\emph{personalizado}), em que o utilizador define tom e
estrutura. Cada \emph{template} combina um \emph{prompt} com metadados de
apresentação.

O desafio linguístico central é que os modelos disponíveis tendem para o
\textbf{português do Brasil}. A resposta foi a constante
\texttt{PT\_PT\_RULES}, embutida em todos os \emph{templates}, com regras
explícitas de léxico, colocação pronominal e proibição de gerúndio
(\Cref{lst:ptpt}). Cada \emph{template} recebe ainda um bloco
\texttt{USER\_FOCUS\_BLOCK} de \emph{prioridade máxima}: a intenção
textual do utilizador determina o \emph{assunto}, enquanto os factos
servem de matéria-prima. O idioma de saída pode ser comutado para inglês
através de um sufixo de \emph{prompt}, sem \emph{templates} paralelos.

\textbf{Rede de segurança determinística.} Como a instrução por si só não
é fiável --- mesmo avisados, o Llama e o Gemini reincidem no português do
Brasil ---, acrescentou-se um \textbf{pós-processador determinístico}
(\texttt{pt\_pt.py}) que corre \emph{depois} da geração e corrige
mecanicamente os brasileirismos mais frequentes: substituições de
vocabulário inequívocas e que preservam género/número (\emph{trem}\,$\rightarrow$\,\emph{comboio},
\emph{fazenda}\,$\rightarrow$\,\emph{quinta}, \emph{caminhão}\,$\rightarrow$\,\emph{camião}\ldots)
e a conversão da construção \emph{auxiliar~+~gerúndio} para
\emph{auxiliar~+~a~+~infinitivo} (``estava chorando''\,$\rightarrow$\,``estava a
chorar''), o sinal gramatical mais reconhecível da variante. O mesmo
módulo expõe uma função \texttt{count\_brasileirismos} usada como métrica
objetiva de aderência ao PT-PT (\Cref{sec:qualidade}). O pós-processador
corrige ainda a \textbf{concordância de género} num conjunto curado de
substantivos (``o~lã''\,$\rightarrow$\,``a~lã''), \textbf{remove parágrafos
repetidos} (que o modelo por vezes produz para encher extensões longas) e
\textbf{corta uma frase final incompleta} deixada por um limite de
\emph{tokens}.

\textbf{Estilo e integração das imagens.} Dois blocos de \emph{prompt}
elevam a qualidade literária: as \texttt{STYLE\_RULES} proíbem a repetição
de nomes próprios (impondo pronomes e perífrases), exigem variedade de
verbos e conectores e reforçam a colocação pronominal europeia; as
\texttt{MEDIA\_GROUNDING} instruem o modelo a situar o espaço, o lugar e as
pessoas \emph{antes} de descrever a cena, com transições suaves, para que
cada fotografia surja em contexto e não de forma abrupta.

\textbf{Controlo de extensão.} O assistente permite escolher a duração em
quatro níveis --- \emph{curta} ($\approx$1\,min), \emph{média}
($\approx$2--3), \emph{longa} ($\approx$4--5) e \emph{épica}
($\approx$6--8). Cada nível gera uma instrução de extensão prioritária (que
sobrepõe o número de parágrafos do \emph{template}) e um limite de
\emph{tokens} calibrado para a duração ($\approx$140~palavras/minuto). Para
não sacrificar a qualidade, a \textbf{não-repetição tem prioridade sobre o
comprimento}: se os factos não bastarem, o modelo escreve um texto mais
curto mas completo.

\subsection{Segmentação em cenas (sincronização narrativa--imagem)}
\label{sub:cenas}

Para que o vídeo deixe de ser um \emph{slideshow} com tempos arbitrários
e passe a um verdadeiro documentário, introduziu-se uma etapa de
\textbf{segmentação em cenas} (\texttt{scene\_builder.py}). Após gerar a
prosa, o M3 divide-a em parágrafos (cada parágrafo $=$ uma cena) e
associa a cada cena as fotografias que melhor a ilustram, produzindo uma
estrutura serializável guardada na coluna \texttt{stories.scenes} --- uma
lista de \texttt{\{text, photo\_ids, caption\}}, cuja forma se ilustra no
\Cref{lst:scene} (Anexo~\ref{ap:codigo}).

A atribuição foto\,$\rightarrow$\,parágrafo é \emph{determinística e sem
dependências externas}: pontua-se a sobreposição léxica entre cada
parágrafo e a ``impressão digital'' textual de cada fotografia (descrição
de IA, \emph{tags}, cenário, sugestão narrativa, \textsc{ocr}, local). É,
na prática, um passo de \emph{retrieval} que ancora os visuais à prosa.
Quando não há sinal léxico, distribuem-se cronologicamente. Numa narrativa
\textbf{longa ilustrada por poucas fotografias}, os parágrafos que sobram
--- onde a narração já não se refere a uma imagem em concreto --- são
preenchidos \textbf{alternando} ciclicamente as disponíveis (sem repetir a
da cena anterior), de modo que o vídeo continua a mudar de imagem em vez de
congelar; cada fotografia surge primeiro no parágrafo que a descreve e só
depois é reaproveitada. Toda a etapa é
defensiva: qualquer falha deixa \texttt{scenes\,=\,None} e o M4 regressa ao
\emph{slideshow} clássico. Esta abordagem responde diretamente ao que o
estudo inicial identificou como o desafio técnico central do M4
(\Cref{sub:geracao_video}): cada fotografia visível exatamente durante a
sua narração.

\subsection{Ancoragem nas fotografias e seleção pelo utilizador}
\label{sub:selecao_fotos}

A narrativa \textbf{não} se baseia apenas na árvore e nas notas que o
utilizador escreve: o construtor de contexto (\texttt{\_build\_events\_from\_media})
incorpora, para cada fotografia, a \textbf{descrição gerada pela visão}, o
\textsc{ocr} (texto de documentos, como certidões), o cenário (o
\emph{espaço} da cena) e o \textbf{lugar} (topónimo derivado do GPS), a
emoção, as \emph{tags} e \emph{quem aparece} na imagem. São estes factos visuais que
ancoram o texto a momentos concretos --- e, por via da segmentação em cenas
descrita acima, determinam também a \textbf{ordem das imagens no vídeo}.

Por defeito, são consideradas todas as fotografias disponíveis (ou, numa
biblioteca grande, as que o \textsc{rag} julga relevantes ao tema). Foi
acrescentada a possibilidade de o utilizador \textbf{escolher explicitamente
quais as fotografias} que entram numa história, diretamente no assistente de
geração: a lista de \texttt{media\_ids} selecionados é enviada ao gerador,
que restringe o contexto a essas imagens (saltando o \emph{retrieval}
automático). Como a narrativa alimenta a segmentação em cenas, esta escolha
propaga-se ao vídeo. O mecanismo funciona tanto no arquivo geral como
\emph{dentro de um projeto} (onde a seleção se limita às fotografias do
projeto), dando ao utilizador controlo curatorial sobre o produto final. Um
projeto pode ainda \textbf{reutilizar fotografias já analisadas} da
biblioteca geral, ligando-as sem novo carregamento nem nova análise --- a
análise de visão é feita uma única vez e reaproveitada ---, para além de
poder receber fotografias carregadas diretamente para si.

A seleção é ainda \textbf{organizada por pessoa}: ao escolher os
intervenientes da história, as fotografias surgem agrupadas por cada pessoa
(as imagens em que está etiquetada, via \texttt{media\_files.person\_ids}),
permitindo escolher exatamente quais das fotos de cada um entram. Em
complemento, cada pessoa tem uma \textbf{galeria própria} onde se anexam ou
removem as suas fotografias e documentos (por exemplo, o retrato e uma
certidão), gerindo o conjunto do lado da pessoa --- o inverso da etiquetagem
``quem aparece'' feita do lado da fotografia.

\subsection{Regeneração orientada por \emph{feedback}}
\label{sub:regenerar}

A geração raramente acerta à primeira no registo pretendido. Para o permitir sem obrigar a recomeçar, o leitor
de histórias oferece uma
\textbf{regeneração \emph{in-place} orientada por \emph{feedback}}: o
utilizador escreve uma nota em texto livre (por exemplo, ``mais curta'',
``foca-te no avô João'' ou ``tom mais sóbrio'') e o sistema reescreve a
narrativa \emph{sobre a mesma história}, reutilizando as mesmas
fotografias, pessoas, idioma e definições de voz/legendas, e preservando o
mesmo identificador (a nota é incorporada no \emph{prompt} como instrução
de prioridade máxima). O mecanismo reusa o gerador do M3 com um parâmetro
\texttt{update\_story\_id}, atualizando o texto e a respetiva segmentação
em cenas em vez de criar um registo novo --- o que evita acumular versões
e mantém o vídeo eventualmente já gerado coerente com a última narrativa.

\section{Módulo M4 --- Geração Multimédia}
\label{sec:impl_m4}

O M4 converte uma narrativa num \textbf{vídeo documental} MP4 (H.264,
$1280\times720$, 24~fps por defeito --- a resolução e os \emph{frames}/segundo
são parametrizáveis por variável de ambiente). O \texttt{M4Processor} orquestra
a seleção cronológica das fotografias, a síntese de voz da narrativa e a
montagem do vídeo, mantendo apenas espaço de trabalho temporário no disco (as
fotografias e o vídeo final vivem no \emph{Storage}). A montagem corre numa
\emph{thread} dedicada e o pedido é tratado de forma \textbf{síncrona}, mantendo
a instância acordada durante o \emph{render}.

\textbf{Geração local, reprodução universal.} A montagem com MoviePy/FFmpeg
é intensiva em memória e a 720p excede os $\approx$512~MB do plano gratuito do
Render (\emph{out-of-memory}, que deixava o vídeo preso em processamento). Em
vez de degradar a qualidade, optou-se por uma política de \textbf{geração
exclusivamente local}: na nuvem, a flag \texttt{VIDEO\_LOCAL\_ONLY} faz o
\emph{backend} \textbf{recusar} de imediato o pedido de \emph{render} (com uma
mensagem clara), ao passo que, em ambiente local, o vídeo é montado a 720p sem
restrição de memória. Como o \emph{backend} local escreve o MP4 no \emph{mesmo}
\emph{Storage} (Supabase) que a nuvem consome, o vídeo \textbf{pré-gerado
localmente} fica disponível e reproduzível no site implantado sem qualquer
\emph{render} em tempo real --- a reprodução é leve (apenas servir o ficheiro) e
nunca é afetada pelo limite de memória. Esta separação entre \emph{produzir}
(local, pesado) e \emph{reproduzir} (nuvem, leve) é uma consequência direta do
requisito de que tudo o que se faz \emph{online} seja também possível
\emph{localmente}, e vice-versa.

\subsection{Montagem cinematográfica}

O \texttt{video\_builder} eleva a qualidade percebida acima de um simples
\emph{slideshow}:

\begin{itemize}[leftmargin=*]
  \item \textbf{Enquadramento não-destrutivo} (\emph{letterbox}): cada
  fotografia cabe inteira no quadro, sem cortes nem distorção, sobre um
  fundo composto pela mesma imagem ampliada, desfocada (\emph{Gaussian
  blur}) e escurecida — um pano de fundo cinematográfico.
  \item \textbf{Movimento por fotografia configurável}
  (\texttt{VIDEO\_MOTION}): \emph{none} (estático), \emph{gentle} ou o
  \textbf{efeito Ken~Burns} --- \emph{zoom} subtil ($1.0\!\rightarrow\!1.06$)
  e \emph{pan} direcional com aceleração suave (\emph{smoothstep}), com o
  ângulo determinístico a partir do \emph{hash} do caminho do ficheiro para
  que fotografias consecutivas não se movam na mesma direção
  (\Cref{lst:kenburns}). Por \textbf{omissão} usa-se o modo \emph{none}
  (fotografias estáticas): além de aliviar a montagem, permite ``pré-cozer''
  uma única vez cada fotograma de cena, acelerando o \emph{render}.
  \item \textbf{Transições} por \emph{crossfade} e \emph{fade-to-black}
  final, evitando cortes abruptos. O \emph{crossfade} usa uma implementação
  própria (\texttt{\_concatenate\_with\_crossfade}) que mistura apenas os
  \emph{frames} sobrepostos, em vez do \texttt{CompositeVideoClip} do
  MoviePy --- uma otimização que reduziu drasticamente o tempo de montagem
  (\Cref{sec:resultados_testes}).
  \item \textbf{Cartões de título e de fim} de abertura e fecho. O restante
  vídeo é mantido \textbf{limpo}, sem texto gravado na imagem: as legendas
  de narração deixaram de ser ``cozidas'' como \emph{lower thirds} e passaram
  a uma \textbf{faixa de legendas comutável} (ver
  \Cref{sub:narracao_legendas}), mais legível e opcional.
  \item \textbf{Ritmo calibrado}: piso de $\sim$4~s por fotografia em modo
  cena (subido de 3~s), para que nenhuma imagem passe a correr contra o
  \emph{crossfade} de 1.6~s.
\end{itemize}

\subsection{Narração, voz e legendas}
\label{sub:narracao_legendas}

A narração é gerada por \textbf{síntese de voz} (vozes neuronais
europeias do \emph{edge-tts}, com \emph{gTTS} como alternativa) na língua
em que a história foi escrita (capturada em \texttt{Story.language}),
garantindo coerência mesmo que o utilizador altere depois o idioma da
interface. O utilizador pode \textbf{escolher a voz} --- masculina ou
feminina --- no assistente de geração (\texttt{Story.voice}); o M4 resolve
essa preferência para a voz neuronal correspondente em cada idioma.

Quando a história tem cenas (\Cref{sub:cenas}), o M4 entra em \textbf{modo
documentário} (\texttt{build\_documentary}): sintetiza \emph{uma narração
por cena} e mostra as fotografias dessa cena durante (aproximadamente) a
duração da sua narração, concatenando os áudios em ordem após o cartão de
título. Assim, o espetador vê a fotografia que está a ser narrada --- a
sincronização semântica que distingue um documentário de uma apresentação.
O \Cref{lst:plandur} mostra o cálculo determinístico (e testado) da duração
de cada fotografia. Na ausência de cenas, usa-se o alinhamento clássico por
duração total. A \Cref{fig:frame_video} (Anexo~\ref{ap:capturas}) mostra um
fotograma do resultado.



\textbf{Legendas como faixa sincronizada (WebVTT).} Em vez de gravar o
texto na imagem, a narração é exportada como uma \textbf{faixa WebVTT}
(\texttt{.vtt}) \emph{comutável} no leitor, sincronizada \textbf{ao nível
da palavra}: o \emph{edge-tts} devolve, com o áudio, os carimbos temporais
de cada palavra (\emph{WordBoundary}), a partir dos quais se escreve o
\texttt{.vtt}. É servido por um \emph{endpoint} próprio e, no leitor, uma
\texttt{<track>} liga/desliga as legendas e escolhe o tamanho (via
\texttt{::cue}). Para o descarregamento, são ainda \emph{embutidas} no MP4
(\emph{mux} \texttt{mov\_text} via FFmpeg). Assim, o vídeo fica limpo por
omissão, mas a acessibilidade textual está sempre disponível.

\section{Infraestrutura do \emph{Backend}}
\label{sec:impl_backend}

\subsection{API e organização}

O \emph{backend} expõe uma API REST versionada (\texttt{/api/v1/})
organizada em grupos de \emph{routers}: \texttt{auth}, \texttt{upload},
\texttt{timeline}, \texttt{narrative}, \texttt{projects},
\texttt{genealogy}, \texttt{multimedia}, \texttt{tasks} e \texttt{health}.
A documentação interativa (\emph{Swagger UI}) é gerada automaticamente.

\subsection{Fila de tarefas e \emph{sweep} de órfãs}

Gerações longas podem correr em segundo plano, com o \emph{frontend} a
acompanhar o estado por \emph{polling} de \texttt{/api/v1/tasks/\{id\}}.
A execução é feita por uma de duas vias, transparentes para o cliente:

\begin{itemize}[leftmargin=*]
  \item \textbf{Celery + Redis} quando existe um \emph{worker} dedicado
  (ambiente completo).
  \item \textbf{Executor \emph{in-process}} quando \texttt{CELERY\_ENABLED}
  é falso (por exemplo, no plano gratuito da nuvem, que não comporta um
  \emph{worker} pago). Neste modo, a tarefa corre num pequeno
  \emph{thread pool} limitado dentro do próprio processo do \emph{backend}
  --- cada \emph{thread} com o seu próprio \emph{event loop}, à imagem de um
  \emph{worker} Celery, para não bloquear o \emph{loop} principal do
  FastAPI. O pedido devolve igualmente um \texttt{task\_id} de imediato.
\end{itemize}

Ambas as vias partilham os mesmos \emph{corpos} de tarefa
(\texttt{tasks/bodies.py}) e o mesmo mecanismo de transição de estado
(\texttt{running}\,$\rightarrow$\,\texttt{done}/\texttt{failed}), pelo que a
experiência do utilizador é idêntica em qualquer ambiente. O \emph{thread
pool} in-process limita-se a um trabalhador --- prudente face aos
$\approx$512~MB do plano gratuito.

Uma decisão de robustez relevante é o \textbf{\emph{sweep} de tarefas
órfãs} no arranque: como nem um \emph{worker} Celery nem um \emph{thread}
\emph{in-process} sobrevivem a um reinício do \emph{backend}, qualquer
tarefa que ficasse \texttt{pending}/\texttt{running} de um arranque
anterior é, por definição, órfã, sendo marcada como falhada --- o que evita
o sintoma de tarefas eternamente ``em execução''.

\subsection{Observabilidade e validação}

O sistema usa \textbf{structlog} para \emph{logging} estruturado (stdout +
ficheiro rotativo). A sonda profunda \texttt{/healthz} verifica todas as
dependências (base de dados, Ollama, Groq, Gemini, Redis, ChromaDB, disco). Os
\emph{uploads} são validados por \emph{magic bytes} (libmagic), tipo MIME
e limites de tamanho (25~MB para fotografias, 100~MB para GEDCOM), e o
\emph{rate limiting} por IP (slowapi) protege contra abuso, com limites
diferenciados (\Cref{tab:rate}).

\begin{table}[H]
\centering
\caption{Limites de taxa aplicados pelo \emph{middleware} (por IP).}
\label{tab:rate}
\small
\begin{tabular}{@{}lc@{}}
\toprule
\textbf{Classe de \emph{endpoint}} & \textbf{Limite} \\
\midrule
Por omissão & 100/minuto \\
Carregamento (\emph{upload}) & 20/minuto \\
Geração (narrativa + vídeo) & 5/minuto \\
\bottomrule
\end{tabular}
\end{table}

\section{\emph{Frontend}}
\label{sec:impl_frontend}

A \textsc{spa} React/TypeScript distingue rigorosamente \emph{estado de
servidor} (gerido pelo \textbf{TanStack Query}, com $\approx$50 \emph{hooks}
que encapsulam pedidos, \emph{cache}, invalidação e \emph{polling}) de
\emph{estado de cliente} (gerido pelo \textbf{Zustand}: autenticação e
tema). Um interceptor do cliente HTTP (axios) injeta automaticamente o
\emph{token} \emph{bearer} em cada pedido e, perante uma resposta
\texttt{401}, termina a sessão de forma limpa. Um \texttt{SessionLoader}
hidrata a sessão Supabase no arranque e instala um \emph{listener} que
mantém o \emph{token} fresco, incluindo renovações silenciosas.

A \textbf{navegação} cobre os nove destinos visíveis da aplicação (Início,
Biblioteca, Família, Linha Temporal, Projetos, Gerar, Histórias, Vídeos e
Definições; a página de \emph{Tarefas} mantém-se acessível pela rota, mas
fora da barra lateral), e cada história exibe um \emph{badge} com o
\textbf{projeto} a que pertence (com ligação direta), clarificando a
organização do arquivo. A
\textbf{página pública} (\emph{landing}) inclui uma \textbf{demonstração
interativa} do \emph{pipeline} (foto $\rightarrow$ análise por IA $\rightarrow$
narrativa $\rightarrow$ vídeo), com a descrição e a história a serem
``escritas'' ao vivo, comunicando o valor do sistema antes do registo.

\textbf{Gestão de histórias.} A lista ganhou ferramentas de curadoria:
\textbf{pesquisa} por título/texto, \textbf{filtro} por projeto,
\textbf{ordenação} e \textbf{favoritos} (estrela que fixa as preferidas no
topo), além de metadados por cartão (fotografias e pessoas usadas). No
leitor, cada história pode ser \textbf{exportada para PDF} e
\textbf{regenerada com \emph{feedback}} (\Cref{sub:regenerar}).

\textbf{Árvore familiar interativa.} A vista \emph{React~Flow}
(\Cref{sec:impl_m2}) permite \textbf{pesquisar e centrar} numa pessoa,
\textbf{destacar a linhagem} ao clicar num nó (o resto da árvore
esbate-se), \textbf{exportar} a árvore como PNG e \textbf{adicionar
familiares} num só passo. Toda a interface está disponível em
\textbf{português europeu} e \textbf{inglês} (\Cref{sec:uiux}), incluindo as
cadeias geradas no cliente.

\subsection{Resiliência a \emph{cold starts}}
\label{sub:coldstart}

O plano gratuito do Render suspende o \emph{backend} ao fim de alguns
minutos de inatividade; o primeiro pedido após a suspensão acorda-o, mas
pode demorar o suficiente para a ligação cair antes de a resposta chegar
ao \emph{browser} — mesmo tendo a operação sido concluída e persistida no
servidor. O \emph{frontend} distingue este caso (resposta perdida, sem
código HTTP) de um erro real e, em vez de mostrar um ``Network Error''
bruto, \emph{reverifica}: aguarda alguns segundos, refaz a consulta
relevante e, se os dados apareceram (uma nova história, novas pessoas
importadas), trata a operação como bem-sucedida. Nunca reenvia o pedido,
para não duplicar dados. Como mitigação de raiz, um \emph{ping} periódico
externo ao \texttt{/healthz} mantém o serviço acordado.
\label{sec:impl_seguranca}

\subsection{Autenticação por JWKS e isolamento multi-inquilino}

A autenticação é delegada ao \textbf{Supabase Auth}. Os tokens são
validados no \emph{backend} contra o \emph{endpoint} JWKS (algoritmo
ES256, chave ECC P-256), com \emph{cache} do documento de chaves em
memória (TTL de 10~min, indexada por \texttt{kid} para suportar rotação de
chaves) e verificação da \emph{audience} (\Cref{lst:jwks}). A identidade
extraída expõe o \textsc{uuid} do utilizador, que é a chave de isolamento:
\textbf{todas} as consultas de domínio filtram por \texttt{user\_id}.

\textbf{Eliminação de conta.} A remoção definitiva de uma conta é feita por um
\emph{endpoint} dedicado que apaga todas as linhas do dono (média, pessoas,
relações, histórias, vídeos, tarefas e projetos) e, em seguida, remove o
utilizador através da \emph{Admin~API} do Supabase (chave \emph{service role}),
garantindo que as credenciais deixam efetivamente de funcionar --- ao contrário
de um mero \emph{sign-out}, que deixava a conta acessível.

\subsection{Servir ficheiros com autenticação}

Os \emph{endpoints} que servem fotografias e vídeos aceitam o \textsc{jwt}
no cabeçalho \texttt{Authorization} \emph{ou} num parâmetro
\texttt{?token=} — necessário porque as \emph{tags} \texttt{<img>} e
\texttt{<video>} não permitem cabeçalhos personalizados. O \emph{backend}
responde com um \emph{redirect} para um URL assinado e temporário do
Supabase Storage, evitando expor as chaves de objeto.

\subsection{Privacidade por inferência local}

A herança \emph{local-first} mantém-se na \textbf{opção} de executar o
\textsc{llm} localmente (Ollama): com o Ollama disponível e sem ligação à
nuvem, a cadeia recai sobre ele e os factos do arquivo não são enviados a
terceiros para a geração de texto. Por omissão, e por razões de desempenho
(\Cref{sec:impl_m3}), o motor de texto na nuvem é o Groq e o Gemini serve a
visão e a rede de segurança final; mas o caminho totalmente local
permanece um \emph{fallback} de primeira classe, garantindo que o sistema
funciona --- e pode ser privado --- mesmo offline.

\section{\emph{Deployment}}
\label{sec:impl_deploy}

A solução é implantada em três serviços geridos: \textbf{Vercel}
(\emph{frontend} estático, com reescrita de rotas desconhecidas para
\texttt{index.html}); \textbf{Render} (\emph{backend} FastAPI a partir de
um contentor \textbf{Docker} próprio, que instala FFmpeg, libmagic e
Tesseract; configuração declarada em \texttt{render.yaml}, com sonda em
\texttt{/healthz}); e \textbf{Supabase} (autenticação, PostgreSQL e
Storage). No ambiente gratuito, o ChromaDB e o \emph{worker} Celery são
desativados — possível precisamente graças às decisões de degradação
graciosa descritas acima. A configuração CORS suporta origens exatas e uma
expressão regular para os \emph{deploys} de pré-visualização do Vercel.
